\section{Introduction}

A gravitational wave is a ripple in the curvature of spacetime. To detect gravitational waves, the Laser Interferometer Gravitational-Wave Observatory (LIGO) relies on the precise measurement of exceedingly small changes in the relative positions of suspended mirrors. They perform this measurement using laser interferometers between mirrors 4~km apart.

Each mirror carries a multilayer optical coating that reflects the laser light. These coatings consist of alternating layers of materials with high and low refractive indices. In Advanced LIGO, titania-doped tantala (Ti:Ta$_2$O$_5$) is used for the high refractive index and silica (SiO$_2$) for the low refractive index. This multilayer structure achieves high reflectivity through constructive interference of light reflected at successive interfaces and maintains the required optical transmission at the same time. However, the coating must satisfy multiple constraints simultaneously. For Advanced LIGO end test masses, the specified transmission at the primary laser wavelength of 1064~nm is $5 \pm 1$ parts per million (ppm). The transmission at 532~nm (used in the interferometer's control system) must remain between 1\% and 4\%~\cite{E0900068}. Furthermore, the coating must minimize thermal noise, particularly coating Brownian noise, as it contributes significantly to the overall displacement noise of the detector~\cite{aLIGO2015,Hong2013}. The mechanical loss of the coating materials introduces an additional constraint in the design process, creating a trade-off between optical performance, mechanical properties, and coating robustness.

Achieving these specifications requires precise control over the thickness of individual coating layers. However, the prescribed thicknesses cannot be produced perfectly. Deviations come from limitations in deposition-rate control, monitoring, and other manufacturing uncertainties~\cite{Tikhonravov2012}. When the light propagates through the coating, even relatively small thickness errors can alter the phase accumulated by it. As a result, it might alter the interference conditions and the resulting optical transmission. In the Advanced LIGO end mirrors, a systematic error in coating-layer thicknesses was reported to leave the transmission at 1064~nm within specification while increasing the transmission at 532~nm by approximately an order of magnitude beyond the desired value~\cite{aLIGO2015}. Thus, manufacturing inaccuracies can compromise the performance of an otherwise carefully optimized coating.

The coating studies we found for gravitational-wave detectors fix the layer design before deposition begins. Venugopalan et al.~\cite{Venugopalan2024} optimized Advanced LIGO-type coatings with the OptimalBragg code and tested the resulting designs against random thickness errors of about 0.5\%. Pierro et al.~\cite{Pierro2021} searched a large set of three-material quarter-wave designs and tested the best ones against independent layer-thickness errors within $\pm 1$~nm and against fluctuations in absorption. In later work, Pierro et al.~\cite{Pierro2026} deposited three-material coatings at LMA and checked the fixed designs against thickness errors of $\pm 0.5$~nm. In each case the design is chosen once, and its tolerance to errors is tested afterwards. Outside the gravitational-wave field, correcting a coating while it is being made is a known technique: the deposited layers are measured, and the layers not yet deposited are re-optimized to compensate~\cite{Sullivan1992,Wilbrandt2008,Melzig2024}. Poitras and Ma~\cite{Poitras2025} simulated this kind of re-optimization after every layer for facet coatings of photonic waveguide devices. We did not find a study that applies it to a gravitational-wave mirror, where the coating has to meet transmission targets at two wavelengths while keeping thermal noise low. Our search was not exhaustive, and coating groups may use similar methods without publishing them.

In this work, we asked whether manufacturing deviations in multilayer mirror coatings can be compensated for by re-optimizing the top four layers, which have not yet been deposited. We used measurements of the already deposited layers to update the thicknesses of the remaining layers. Our objective was to recover the required optical transmission while maintaining the coating's thermal-noise performance. We refer to this procedure as coating rescue. We tested and examined a population of 32 distinct coating designs that initially satisfied the relevant optical specifications.

We used Monte Carlo simulations to investigate the effectiveness of this approach under both independent layer-to-layer deposition errors and systematic errors shared across the coating. Then we also examined how measurement uncertainty affected the effectiveness of re-optimization and evaluated the resulting changes in coating Brownian noise. To be specific, we quantified the deposition-error tolerance using the error level at which 95\% of simulated coatings continued to satisfy both optical transmission requirements. Our three questions were: (1) how large the deposition errors can be while 95\% of mirrors still pass, (2) how this depends on measurement error, and (3) whether the rescue increases coating Brownian noise. Through these investigations, we aim to establish the conditions under which measurement-informed re-optimization can improve the manufacturing tolerance of gravitational-wave detector mirror coatings, as well as the limitations imposed by measurement accuracy and systematic deposition errors.
